\section{The complete twelve-root extraction}
\label{sec:explicit-twelve-roots}

This section gives the complete polynomial identity behind the finite
testing family. All lower-case source parameters in this section are in
the coefficient field. The variables $u,w,W,X,Y,Z$ are independent
testing variables. Put $q=\sqrt z$, $x=k^2$, $Q_b=1-b$, and use the
native stationary C1 source of (P1). In particular
\[
 a=p/b,\quad \ell=l/Q_b,\quad e=xh/\sigma,\quad
 \rho=\lambda x/\sigma,\quad L_0=\lambda\sigma,\quad
 q_0=\lambda k,\quad d_G=4D/\sigma^2,\quad S=d_G+\rho.
\]
Here $\lambda>0$. Set $a_0=k/\sigma$, $c=\max(0,1-K_A)$ and
$K_+=K_A+c\ge1$. No inverse of $D$, $\beta$, or $c$ is used.

The polynomial is the complete parent (P4), with all four state
coordinates independent. Its full source definitions are (P1) and (SP3).


\subsection{The twelve principal roots}

Define
\[
\begin{gathered}
 N_x=X+uhZ,\qquad N_y=(Y-bwZ)/q,\qquad R=wZ/q-uqW,\\
 \Gamma=\sigma Z-uwW,\qquad G_* =\sqrt{d_G}\,\Gamma,\\
 A_\theta=q(u^2-k^2)W,\qquad T_\theta=kq(u-k)W,\\
 U_* =\sqrt{q_0}\,q\sqrt u(u-k)W,\qquad
 K_* =\sqrt{q_0k}\,q(u-k)W.
\end{gathered}
\]
For $u,w,W,Z\ge0$ put
\[
\begin{gathered}
 B_1=\sqrt{L_0}\,uZ,\quad B_2=\sqrt\rho\,uwW,\quad
 B_3=\sqrt\lambda\,k^2qW,\quad B_4=\sqrt\lambda\,kwZ/q,\\
 G=\sqrt\lambda\,k\sqrt{uwWZ},\qquad d_i=B_i-G.
\end{gathered}
\]
Thus
\[
 \psi=(N_x,N_y,R,G_*,A_\theta,T_\theta,d_1,d_2,d_3,d_4,U_*,K_*)^{\mathsf T}
       =A_{\rm src}\phi,
\]
where $A_{\rm src}$ has only source entries and
\[
 \phi=(X,Y,Z,W,uZ,wZ,uW,uwW,u^2W,u^{3/2}W,u^{1/2}W,
             \sqrt{uwWZ})^{\mathsf T}.
\]
This proves the claimed separation-rank bound twelve. In particular,
the fourth-root circuit does not add an independent principal testing
function: $B_1B_2B_3B_4=\lambda^2x^2u^2w^2W^2Z^2$ gives the printed
formula for $G$ exactly.

If $W\ge0$ and $Z\le0$, retain the same signed $B_i$ and put $d_i=B_i$.
The twelfth column of $A_{\rm src}$ is then zero. The full matrix is
printed in \texttt{C1\_ROOT\_MATRIX.json} for both signs.

\subsection{Both state-normal conductors and the exact Gram}

Write $J_C=N_y+bR$ and $J_\ell=N_y-Q_bR$. The two linear forms
\begin{align*}
 \mathfrak U={}&\frac{\sqrt{d_G}uG_*
 +\sqrt\rho\{k(d_1-d_2)-(u-k)(d_2-d_3)\}
 -\rho\sqrt\sigma T_\theta}{S},\\
 \mathfrak V_0={}&\frac{R_0\sqrt{d_G}qwG_*+R_0\rho\sigma q^2R
 -R_0\sqrt\rho(qw+q^2\sqrt\sigma)(d_2-d_3)}{S}\\
 &+\frac{R_0q_0q^2+(R_0\sqrt\lambda)\sqrt\rho qw}{S}T_\theta,
 \qquad \mathfrak V=\mathfrak V_0-8pqwu\mathfrak U
\end{align*}
satisfy, exactly,
\[
 \mathfrak U=u\Gamma,\qquad
 \mathfrak V=(R_0-8pu^2)qw\Gamma.
\]
Both complete normal terms are therefore retained. These formulas
are also literal when $D=0$.

Let
\[
 E_{RR}^0=z\{8bD+8bp\sigma+4b\sigma(2C+2D+1)(r-1)\},
\]
and let $E_{yR},\nu,R_0,R_x,B_x,A_y,\chi$ have the explicit source
definitions in (SP3). Put
\[
 \varepsilon_C=2/\lambda+u+k,\qquad
 \varepsilon_G=2a_0u+(u+k)/L_0.
\]
The exact principal Gram is the following quadratic form in the
twelve formal coordinates $\psi$:
\begin{align*}
 \psi^{\mathsf T}\mathcal Q\psi={}&
 2N_x^2+J_C^2+J_\ell^2+G_*^2
 +\frac{K_+}{2}(A_\theta^2+T_\theta^2)
 +\frac12\sum_{i=1}^4d_i^2+U_*^2+\frac12K_*^2\\
 &+\left(E_{RR}^0-\frac{\lambda x}{2}-\frac{cd_Gu^2}{S}\right)R^2
 +E_{yR}N_yR-4\nu N_x\mathfrak U+J_C\mathfrak V\\
 &-4uqw(2m+2D+1)N_xN_y+uqwR_xN_xR
 +qwB_xN_xA_\theta+uA_yN_yA_\theta+\chi uA_\theta R\\
 &+\sqrt{q_0u}\,U_*R
 -\frac cS(\sqrt{a_0u}\,U_*+\sqrt{a_0k}\,K_*)^2
 -\frac{c\varepsilon_G}{S}G_*^2
 -\frac{cd_G\varepsilon_C}{S}\sum_{i=1}^4d_i^2.
\end{align*}
In $s=\sqrt u,w$, this matrix has just the testing monomials
\[
 1,\quad w,\quad w^2,\quad s,\quad s^2,\quad s^2w,
 \quad s^4,\quad s^4w.
\]
The eight evaluated symmetric coefficient matrices, comprising 104
stored upper-triangle entries, are in
\texttt{C1\_PRINCIPAL\_GRAM.json}.

\subsection{The nineteen elementary retained roots}

Define
\[
 X_e=X+ewW,\qquad F_A=u(Z-a_0wW),
\]
\[
 k_1=2(1+2m)(h-Q_b)^2+8Dhb+8pb\sigma+4\beta(1-h)^2.
\]
The following nineteen roots are retained without any mixed square
root. Every displayed radical is purely source-side.
\begin{align*}
 &2\sqrt\beta N_x,\quad2\sqrt\beta wN_x,
 \quad\sqrt{2(2D+1)}uX_e,\quad2\sqrt{bp}\,qu\Gamma,\\
 &2\sqrt a\,uJ_C,\quad2\sqrt a\,wJ_C,
 \quad2\sqrt a\,uwJ_C,\quad2\sqrt\ell\,uJ_\ell,
 \quad2\sqrt C\,qwJ_C,\\
 &\sqrt{2+4m}\,wN_x,\quad
 \sqrt{4a+2z(2C+1)-1}\,J_C,\quad\sqrt{4\ell-1}\,J_\ell,\\
 &\sqrt{2(1+2m)}\,quN_y,\quad
 \sqrt{2(1+2D)}\,qwN_y,\quad
 \sqrt{k_1}\,quR,\quad\sqrt{2(2D+1)}\,qbwR,\\
 &\sqrt{L_0/2}\,F_A,\quad
 \sqrt{K_+/2}\,qu(u-k)W,\quad
 \sqrt{K_+k}\,q\sqrt u(u-k)W.
\end{align*}
The final pair represents $K_+(A_\theta^2-T_\theta^2)/2$.
The native source bounds give $4a>1$ and $4\ell>1$, so their
radicands are nonnegative. Each root is a source-linear combination
of at most three fixed testing monomials.

\subsection{Eight angular roots, and the negative-state correction}

Let
\[
 \mathscr C=\sum_{i=1}^4B_i^2-4\lambda xuwWZ,
 \qquad
 \Delta_A=\varepsilon_C\mathscr C+u^2R^2
                   +\varepsilon_G\Gamma^2-A_\theta^2.
\]
For $W,Z\ge0$ put
\[
 L=\sqrt{q_0\sigma}\,F_A+\frac{k}{\sqrt{q_0\sigma}}\Gamma,
 \quad K_Z=k^2qW-kwZ/q,\quad H_Z=\sqrt{wWZ}.
\]
The elementary angular identity used here is, for $U,V,Z_*\ge0$
and $a_*>0$,
\[
 \Psi=1+U^2V^2+U^2Z_*^2+V^2Z_*^2-4UVZ_*,
 \quad T_*=U^2-1,\quad Q_*=VZ_*-U,
\]
\[
 R_*=Z_*-UV,\qquad G_*=U(V-Z_*),
\]
\begin{align*}
 &[2+a_*(U+1)]\Psi+U^2Q_*^2+
       [2U+(U+1)/a_*]R_*^2-T_*^2\\
 &=(T_*+UQ_*)^2+2Q_*^2
       +(U+1)(\sqrt{a_*}G_*-R_*/\sqrt{a_*})^2\\
 &\qquad+a_*(U+1)(1-VZ_*)^2+
                  2a_*(U+1)VZ_*(U-1)^2 .
 \tag{US}
\end{align*}
Multiplication verifies it. For $W\ne0$, substitute
\[
 U=u/k,\quad V=w/(\sqrt{\sigma}q),\quad
 Z_*=\sqrt{\sigma}Z/(kqW),\quad a_*=q_0
\]
and multiply by $(k^2qW)^2$. The resulting formulas below no longer
divide by $W$ or $Z$, so their polynomial identity includes those
zero coordinates.

Then $\Delta_A$ is the sum of the squares of exactly these eight roots:
\begin{align*}
 &A_\theta+uR,\qquad\sqrt2kR,\qquad
 L,\qquad\sqrt{u/k}\,L,\\
 &\sqrt{q_0}K_Z,\qquad\sqrt{q_0u/k}K_Z,
 \qquad\sqrt{2q_0k}(u-k)H_Z,
 \qquad\sqrt{2q_0u}(u-k)H_Z.
\end{align*}
This follows by multiplying out (US);
the displayed formulas no longer divide by a state coordinate.
Multiply these roots by the pure-source factor $\sqrt{cd_G/S}$.

If $W\ge0,Z\le0$, first use the eight formulas with $|Z|$ in place
of $Z$, and append the three roots
\[
 2k\sqrt{4+2q_0+q_0^{-1}}\sqrt u H_{|Z|},\quad
 2\sqrt k\sqrt{2q_0+2+q_0^{-1}}uH_{|Z|},\quad
 2u\sqrt u H_{|Z|}.
\]
Their squared sum is exactly the increase in $\Delta_A$ when the
positive state coordinate is replaced by its negative. This proves
the eleven-root formula on the negative-state chart.

\subsection{The three circuit roots and their fixed testing span}

On the positive-state chart set
\[
 a_t=u^{1/4},\quad b_t=w^{1/4},\quad
 \omega=W^{1/4},\quad\zeta=Z^{1/4},\quad
 \gamma=\sqrt{\sqrt\lambda\,k}.
\]
The six testing monomials are
\[
\begin{gathered}
 f_1=a_t^3b_t\omega\zeta^3,\quad
 f_2=a_t^3b_t^3\omega^3\zeta,\quad
 f_3=a_tb_t\omega^3\zeta,\quad
 f_4=a_tb_t^3\omega\zeta^3,\\
 f_5=a_t^3b_t^2\omega^2\zeta^2,\qquad
 f_6=a_tb_t^2\omega^2\zeta^2.
\end{gathered}
\]
The retained $\mathscr C_+=\mathscr C-\sum d_i^2$ is the sum of
the squares of
\begin{align*}
 &\sqrt2\gamma(L_0^{1/4}f_1-\rho^{1/4}f_2),\\
 &\sqrt2\gamma\{(\sqrt\lambda k^2q)^{1/2}f_3
                    -(\sqrt\lambda k/q)^{1/2}f_4\},\\
 &2\gamma\{(\lambda k)^{1/4}f_5-(\lambda k^3)^{1/4}f_6\}.
\end{align*}
On the negative-state chart there is instead the single root
$2\sqrt\lambda k\sqrt{uwW|Z|}$.

The common scalar coefficient retained in the parent is
\[
 q_C=\frac12-\frac{cd_G}{S}(2/\lambda+k+u).
\]
On $0\le u\le\delta^2$, with $u=\delta^2A_t^8$ in the scaled
testing parametrization, write $A_C=cd_G\delta^2/S$ and
$q_C^0=\frac12-cd_G(2/\lambda+k)/S$. Exactly
\[
 q_C=(q_C^0-A_C)+A_C(1-A_t^8),\qquad
 1-A_t^8=\xi^2(1+A_t^4),\quad\xi^2=1-A_t^4.
\]
Thus the three fixed multipliers $1,\xi,\xi A_t^2$ suffice for
each circuit root. There is no point-dependent choice of a radical
or a bounding constant.

\subsection{Exact assembly and boundary specialization}

The principal quadratic form, the nineteen elementary squares,
the weighted angular squares, and $q_C\mathscr C_+$ sum to $F_w$.
This identity was independently verified first against the complete
SP6 expression and then against the original P4 parent after the
generic rational P1 substitution. In the latter check all seventeen
nonzero pre-substitution coefficients in the independent polynomial
variables $W,X,Y,Z,u,w$ vanish identically.

All final testing functions are polynomial in quarter-angle and
quarter-state coordinates. Consequently $u=0,w=0,W=0,Z=0$ are
literal specializations, with no limiting SOS argument. The signs
of the quarter-state parameters do not matter: each circuit root
has one fixed parity in each parameter, so its square is invariant.
The testing sign gauges
$(u,Y,Z)\mapsto(-u,-Y,-Z)$ and
$(w,X,Z)\mapsto(-w,-X,-Z)$ restore the other angle signs. An overall
state sign makes $W\ge0$; the two state charts then cover all real
states. Real and imaginary parts restore complex states.

The complete retained coefficient matrices are in
\texttt{C1\_RETAINED\_ROOTS.json}: the positive-state chart has
33 distinct retained testing monomials and 30 roots before the
circuit coefficient is split; the negative-state chart has 27
monomials and 31 roots. Every retained root has separation rank at
most three. These are local certificates for the printed source
and testing domains; none of these identities asserts that this
one chart is an exhaustive global Q222 testing family.
